\documentclass[11pt]{article}
\usepackage[margin=1in]{geometry}
\usepackage{amsmath,amssymb,amsthm}
\usepackage{microtype}
\usepackage{booktabs}
\usepackage[hidelinks,pdfusetitle]{hyperref}

\numberwithin{equation}{section}
\newtheorem{theorem}{Theorem}[section]
\newtheorem{lemma}[theorem]{Lemma}
\newtheorem{proposition}[theorem]{Proposition}
\newtheorem{corollary}[theorem]{Corollary}
\DeclareMathOperator{\arcosh}{arcosh}
\newcommand{\R}{\mathbb R}
\newcommand{\C}{\mathbb C}
\newcommand{\ee}{\mathrm e}
\newcommand{\ii}{\mathrm i}
\newcommand{\dd}{\,\mathrm d}
\newcommand{\ip}[2]{\langle #1,#2\rangle}

\title{The Fixed Jacobi Chain for Newman Flow\\
\large Spectral realization and finite-chain readout bounds}
\author{Edward Baker}
\date{14 November 2026}

\begin{document}
\maketitle

\begin{abstract}
The complete theta source at Newman time \(1/2\) defines a probability
measure whose orthogonal polynomials realize multiplication as a
self-adjoint Jacobi operator. Height and reversed Newman time act on
one prepared state by the commuting multipliers \(\exp(\ii xJ_*)\)
and \(\exp(-\tau J_*^2)\). The state admits an exact bond-two encoding
on the one-excitation sector. Cutting the Jacobi matrix at \(N\) sites
gives the \(N\)-node Gaussian quadrature for the same source. A Hermite
remainder averaged over a real Gaussian bounds the normalized
absolute scalar error by the monic norm times an explicit polynomial
in \(\tau\) and \(x\). The complete-source score bounds the monic norm
and yields a sufficient depth of order \(K\log(K+1)\), where
\(K=X+\log(1/\epsilon)+1\) for heights \(|x|\le X\) and normalized
absolute tolerance \(\epsilon\). We also derive an uncertainty bound
for the finite matrix and illustrate it with an existing 16-site
interval certificate.
\end{abstract}

\section{The complete theta source}\label{sec:source}

The finite-chain approximation concerns a scalar observation of a
fixed source. We first fix its physical normalization. Write
\[
 \vartheta(t)=\sum_{m\in\mathbb Z}\ee^{-\pi m^2t},\qquad
 p(u)=\frac{\ee^u}{2}\bigl(\vartheta(\ee^{4u})-1\bigr),\qquad
 r_m(u)=\pi m^2\ee^{4u}.
\]
The Jacobi transformation \(\vartheta(t)=t^{-1/2}\vartheta(1/t)\)
for \(t>0\) \cite{theta} gives
\[
 p(-u)-p(u)=\sinh u,\qquad p'(0)=-\frac12.
\]
Define the amplitude and field by
\begin{align}
 \Phi(u)&=\frac{p''(u)-p(u)}8
     =\ee^u\sum_{m\ge1}(2r_m(u)^2-3r_m(u))\ee^{-r_m(u)},
       \label{eq:source}\\
 W_\lambda(u)&=\ee^{\lambda u^2}\Phi(u),\qquad
 H_\lambda(z)=\frac12\int_{\R}W_\lambda(u)\ee^{\ii zu}\dd u.
       \label{eq:field}
\end{align}
The series is locally differentiable at every order. Applying
\(\partial_u^2-1\) to the reflection relation proves that \(\Phi\)
is even. Every summand is positive on \(u\ge0\), where \(r_m\ge\pi\);
hence the complete \(\Phi\) is positive on \(\R\).

\begin{lemma}[Domination]\label{lem:domination}
Set \(\rho(u)=\pi\ee^{4|u|}\). For each integer \(j\ge0\) there is a
finite constant \(C_j\) such that
\[
 |\Phi^{(j)}(u)|\le C_j \ee^{|u|}\rho(u)^{j+2}\ee^{-\rho(u)} .
\]
In particular, \(|\Phi^{(j)}(u)|\ee^{L u^2+b|u|}(1+|u|)^k\)
is integrable and tends to zero at both ends for every fixed
\(L,b,k\ge0\).
\end{lemma}
\begin{proof}
On the positive half-line, the \(j\)-th differentiated summand is
\(\ee^u Q_j(r_m)\ee^{-r_m}\), where \(Q_0(r)=2r^2-3r\) and
\(Q_{j+1}=(1-4r)Q_j+4rQ_j'\). Its degree is \(j+2\).
Since \(r_m=m^2\rho\) there, the absolute summand is bounded by
\(\ee^u\rho^{j+2}\ee^{-\rho}\) times a constant depending on \(j\)
times \(m^{2j+4}\ee^{-\pi(m^2-1)}\). This last sequence is summable.
The series and its derivatives are also locally uniformly convergent
on any compact interval of the full line, where \(r_m\ge c m^2\)
for some \(c>0\). Evenness therefore extends the estimates with
absolute values to the full line. After taking logarithms, all
positive terms in the last claimed bound grow at most quadratically
in \(|u|\), while \(\rho(u)\) grows exponentially. The resulting
tail is bounded by a constant times \(\ee^{-\rho(u)/2}\).
\end{proof}

Evenness changes \eqref{eq:field} into a half-line cosine integral.
To check the normalization, let
\(I(z)=\int_0^\infty p(u)\cos(zu)\dd u\).
The half-line series for \(p\) and \(p'\) have the same
double-exponential domination used in Lemma~\ref{lem:domination}.
Integration by parts twice gives
\begin{equation}\label{eq:primitive}
 8H_0(z)=\frac12-(1+z^2)I(z).
\end{equation}
The constant arises from \(-p'(0)\). For \(\Re s>1\), integrating
the series for \((\vartheta(t)-1)/2\) against \(t^{s/2-1}\) gives
\(\pi^{-s/2}\Gamma(s/2)\zeta(s)\) by the gamma integral.
Split at \(t=1\), apply the Jacobi transformation on \((0,1)\), and
substitute \(t=1/v\) in the reflected term. The elementary terms
integrate to \(1/(s-1)-1/s=1/[s(s-1)]\). Consequently
\[
 \xi(s)=\frac12+\frac{s(s-1)}2
   \int_1^\infty\frac{\vartheta(t)-1}{2}
     \left(t^{s/2-1}+t^{(1-s)/2-1}\right)\dd t,\qquad
 \xi(s)=\frac{s(s-1)}2\pi^{-s/2}\Gamma(s/2)\zeta(s).
\]
The integral is entire in \(s\) by exponential decay at infinity.
Under \(t=\ee^{4u}\) and \(s=(1-\ii z)/2\) it becomes \(8I(z)\).
Using \(4s(s-1)=-(1+z^2)\) in \eqref{eq:primitive} proves
\begin{equation}\label{eq:xi}
 H_0(z)=\frac18\,\xi\!\left(\frac12-\frac{\ii z}{2}\right).
\end{equation}
Lemma~\ref{lem:domination} also allows differentiation under the
integral on compact sets in \((\lambda,z)\in\C^2\).
In particular \(\partial_\lambda H_\lambda=-\partial_z^2H_\lambda\).
We work on \(0\le\lambda\le1/2\) and put \(\tau=1/2-\lambda\).

\section{The fixed Jacobi operator and its orbit}\label{sec:operator}

Set
\[
 W_*=W_{1/2},\qquad Z_*=\int_{\R}W_*(u)\dd u=2H_{1/2}(0),
 \qquad \dd\mu_*(u)=Z_*^{-1}W_*(u)\dd u.
\]
Use the complex \(L^2(\mu_*)\) inner product linear in the second
slot. Let \(X_{\rm mul}\) denote multiplication by \(u\) on
\[
 \mathcal D(X_{\rm mul})=\{f\in L^2(\mu_*):uf\in L^2(\mu_*)\}.
\]
This operator is self-adjoint. It is densely defined by bounded
truncation, symmetric because \(u\) is real, and
\(X_{\rm mul}\pm\ii\) are onto with bounded inverse multiplication by
\(1/(u\pm\ii)\). Its spectral projections are multiplication by
indicators of real Borel sets.

\begin{lemma}[Polynomial basis and core]\label{lem:core}
The real orthonormal polynomials \(p_n\) for \(\mu_*\), with positive
leading coefficients and \(p_0=1\), form an orthonormal basis. Their
finite span is a core for \(X_{\rm mul}\).
\end{lemma}
\begin{proof}
All finite Gram matrices are positive definite by strict positivity
of the density on \(\R\). To see the needed density, take either
\(\nu=\mu_*\) or \(\nu=(1+u^2)\mu_*\). If \(g\in L^2(\nu)\) is
orthogonal to all powers of \(u\), then
\(F(z)=\int g(u)\ee^{\ii zu}\dd\nu(u)\) is entire by
Cauchy--Schwarz and Lemma~\ref{lem:domination}, and all its Taylor
coefficients at zero are zero. The measure \(g\dd\nu\) has an
\(L^1(\R)\) density. Fourier uniqueness gives that this density
vanishes. For this use, uniqueness follows by convolving with
Gaussian densities. The Gaussian Fourier formula and Fubini give
zero for every such convolution, and the Gaussians form an
approximate identity in \(L^1\). Since \(\nu\) has positive density,
\(g=0\). The second choice of \(\nu\) gives density in the graph
norm \((\|f\|_2^2+\|X_{\rm mul}f\|_2^2)^{1/2}\), proving the core
claim.
\end{proof}

Let \(P_n\) be the monic polynomials at the base slice, and write
\(h_n^*=\int P_n^2\dd\mu_*\), so \(h_0^*=1\).
Evenness gives \(P_n(-u)=(-1)^nP_n(u)\) by uniqueness.
In the expansion of \(uP_n\) in the monic basis, the leading
coefficient is one, the \(P_n\) coefficient vanishes by parity, and
terms of degrees \(k\le n-2\) vanish since \(uP_k\) has degree at most
\(n-1\). The coefficient of \(P_{n-1}\) is \(h_n^*/h_{n-1}^*\),
by moving \(u\) onto \(P_{n-1}\) in the inner product. Thus
\[
 up_n=a_{n+1}p_{n+1}+a_np_{n-1},\qquad
 a_n=\sqrt{h_n^*/h_{n-1}^*}>0\ (n\ge1),\qquad a_0=0.
\]
The \(n=0\) formula omits the second term. Under the unitary map
\(\mathcal U:\ell^2(\mathbb N_0)\to L^2(\mu_*)\) with
\(\mathcal U e_n=p_n\), define
\[
 J_*=\mathcal U^{-1}X_{\rm mul}\mathcal U,\qquad
 \mathcal D(J_*)=\{v:\ u\mathcal Uv\in L^2(\mu_*)\}.
\]
On finite sequences \(J_*\) is tridiagonal, with zero diagonal and
edge \(a_n\) between sites \(n-1\) and \(n\). The core proves that
the closure of this finite-sequence action is the displayed
self-adjoint operator. More generally the domain of \(J_*^m\) consists
of \(v\) with \(u^m\mathcal Uv\in L^2(\mu_*)\). Orthogonal-polynomial
chain mappings are familiar in reservoir theory \cite{chin}.
Here the full-line density, which has unbounded support, determines
the realization by Lemma~\ref{lem:core}.

\begin{theorem}[Height and reversed time]\label{thm:orbit}
For real \(x\) and \(0\le\tau\le1/2\), define
\[
 v(\tau,x)=\ee^{-\tau J_*^2}\ee^{\ii xJ_*}e_0 .
\]
Then \(H_{1/2-\tau}(x)=(Z_*/2)\ip{e_0}{v(\tau,x)}\).
On the prepared state, \(\partial_xv=\ii J_*v\) and
\(\partial_\tau v=-J_*^2v\). The corresponding readout extends to
complex \(z\) by using the prepared vector for the multiplier
\(\ee^{-\tau u^2+\ii zu}\).
\end{theorem}
\begin{proof}
The transform of \(e_0\) is \(1\), so the spectral representation
gives \(\mathcal Uv=\ee^{-\tau u^2+\ii xu}\), which proves the
readout by \eqref{eq:field}. For every fixed \(m\), the same expression
with \(z\in\C\) satisfies \(u^m\mathcal Uv\in L^2(\mu_*)\), with
local domination from Lemma~\ref{lem:domination}. This also
justifies the two strong derivatives, including at the time
endpoints by continuation of the prepared vector. For real heights,
the unitary group \(\ee^{\ii xJ_*}\) and the contraction semigroup
\(\ee^{-\tau J_*^2}\) commute. For nonreal heights, a spectral
multiplier with value \(g(u)\) has domain
\(\{v:g\mathcal Uv\in L^2(\mu_*)\}\); in particular the bare
\(\ee^{\ii zJ_*}\) is unbounded for \(\Im z\ne0\). The formula at
such \(z\) uses the stated vector in its domain.
\end{proof}

In components the two generators give
\begin{align}
 \partial_x v_n
    &=\ii(a_nv_{n-1}+a_{n+1}v_{n+1}),\notag\\
 \partial_\tau v_n
    &=-a_na_{n-1}v_{n-2}-(a_n^2+a_{n+1}^2)v_n
       -a_{n+1}a_{n+2}v_{n+2},                       \label{eq:lattice}
\end{align}
with terms at negative indices omitted. The identities follow
by testing against \(J_*e_n\) and \(J_*^2e_n\), which are finitely
supported. The height generator has range one and the time
generator has range two in the site index. A finite height
increment uses the full unitary group. The chain is fixed at its
base slice. Rebuilding the orthogonal basis for each \(W_\lambda\)
instead gives a family of moving recurrence coefficients.

\section{A single-excitation representation}\label{sec:mps}

On sites \(0,1,\ldots\) take the infinite tensor product of
\(\C^2\) with reference vector \(|0\rangle\). Concretely, complete the
span of mutually orthogonal occupation configurations with
finitely many \(1\)'s and unit norm. The one-excitation sector is
the closed span of \(|1_n\rangle\), with occupation only at site
\(n\). The map \(v\mapsto\sum_n v_n|1_n\rangle\) is an isometry.
For the first \(N\) sites, a matrix product state of bond dimension
two uses
\[
 A_n^0=I_2,\quad A_n^1=v_nE_{12},\quad
 \ell=(1,0),\quad r=(0,1)^{\mathsf T},
 \qquad E_{12}=\begin{pmatrix}0&1\\0&0\end{pmatrix}.
\]
For an occupation word, the amplitude is the corresponding core
product between \(\ell\) and \(r\). The vacuum amplitude is zero,
a single excitation at \(n\) has amplitude \(v_n\), and all words
with at least two excitations vanish since \(E_{12}^2=0\).
After filling the later sites with \(|0\rangle\), the finite
states converge in norm to the full state, with squared error
\(\sum_{n\ge N}|v_n|^2\). This encoding works for every
\(\ell^2\) vector. It stores the amplitudes in infinitely many cores
and its convergence is a statement about these state vectors.

For fixed \(\tau\ge0\), \(v_n(\tau,x)\) is the Fourier transform of
the \(L^1\) density \(p_n(u)\ee^{-\tau u^2}W_*(u)/Z_*\). It tends
to zero as \(|x|\to\infty\): compactly supported step functions
are dense in \(L^1\), and the integral of \(\ee^{\ii xu}\) over
any fixed interval tends to zero. Yet
\[
 \|v(\tau,x)\|_2^2=\int_\R\ee^{-2\tau u^2}\dd\mu_*(u)>0
\]
is constant in real height. Every fixed prefix therefore loses
all its mass as \(|x|\to\infty\). In this basis, a fixed finite
prefix cannot approximate the state uniformly over the whole
real height axis.

\section{Finite principal chains and the scalar remainder}\label{sec:quad}

Let \(\Pi_N\) project onto \(e_0,\ldots,e_{N-1}\) and take
\(J_N=\Pi_NJ_*\Pi_N\) on that range. The finite dynamics uses
\(J_N^2\). In terms of the standard rank-one projection onto the
last site,
\begin{equation}\label{eq:boundary}
 \Pi_NJ_*^2\Pi_N\big|_{\operatorname{ran}\Pi_N}
       =J_N^2+a_N^2\, e_{N-1}e_{N-1}^* .
\end{equation}
Indeed \((1-\Pi_N)J_*\Pi_N\) has only the map from \(e_{N-1}\)
to \(a_Ne_N\). The finite observation is
\begin{equation}\label{eq:finite}
 H_N(\tau,x)=\frac{Z_*}{2}
      \ip{e_0}{\ee^{-\tau J_N^2}\ee^{\ii xJ_N}e_0}.
\end{equation}
Projecting the exact evolved state would retain its \(e_0\)
coordinate unchanged. Formula \eqref{eq:finite} instead evolves
the original \(e_0\) by a new finite generator.

\begin{proposition}[The literal Gaussian rule]\label{prop:gauss}
Let \(Q_N\) be the spectral measure of \(J_N\) at \(e_0\).
It is an even probability measure on the \(N\) simple zeros of
\(P_N\), with positive weights. It agrees with \(\mu_*\) through
degree \(2N-1\), and its first missed moment is
\[
 \int u^{2N}\dd\mu_*-\int u^{2N}\dd Q_N=h_N^*
            =\prod_{j=1}^N a_j^2 .
\]
In particular \(H_N\) is exactly the \(N\)-node Gaussian
quadrature of the complete-base integral.
\end{proposition}
\begin{proof}
The tridiagonal determinant recursion is the monic recurrence
for \(P_N\), including the initial polynomials \(1,u\). Any
eigenvector of \(J_N\) is determined by its first component by
forward recursion, and that component is nonzero for a nonzero
eigenvector. The finite real symmetric matrix has simple
eigenvalues, and its spectral weights are squares of the first
components of normalized eigenvectors. The parity operator
\(e_j\mapsto(-1)^je_j\) conjugates \(J_N\) to \(-J_N\) and
fixes \(e_0\), giving evenness.

For any integer \(k\ge0\), the two moments are
\(\ip{e_0}{J_*^ke_0}\) and \(\ip{e_0}{J_N^ke_0}\),
respectively. Each matrix element is a sum over length-\(k\)
walks along the chain returning to site \(0\). A walk reaching
site \(N\) takes at least \(2N\) steps to return. At length
\(2N\) there is exactly one such walk, with weight
\(\prod_{j=1}^Na_j^2=h_N^*\). This proves exactness and identifies
the first defect. The finite spectral expression for
\eqref{eq:finite} integrates \(\ee^{-\tau u^2+\ii xu}\) against
\(Q_N\) with the factor \(Z_*/2\). These are the classical
Gaussian nodes, weights and degree of exactness \cite{gauss}.
\end{proof}

\begin{theorem}[Gaussian-averaged Hermite remainder]\label{thm:hermite}
For \(N\ge1\), \(x\in\R\), and \(0\le\tau\le1/2\),
\begin{equation}\label{eq:hermite}
 \frac{|H_{1/2-\tau}(x)-H_N(\tau,x)|}{Z_*/2}
   \le h_N^* S_N(\tau,x),\qquad
 S_N(\tau,x)=\sum_{k=0}^N\frac{\tau^k x^{2N-2k}}{k!(2N-2k)!}.
\end{equation}
\end{theorem}
\begin{proof}
For \(s\in\R\) let \(L_s\) be the polynomial of degree at most
\(2N-1\) that matches \(\cos(su)\) and its first derivative at
the nodes. Such a polynomial is unique, because any difference
is divisible by \(P_N^2\); the square linear interpolation
system is therefore invertible. At a point \(u\) outside the
nodes, subtract \(L_s\) and a multiple of \(P_N^2\) from the
cosine so the difference also vanishes at \(u\). Counting the
double zeros and applying repeated real Rolle on the finite
interval containing \(u\) and the nodes identifies the multiple
as the \(2N\)-th derivative of the cosine at a point divided by
\((2N)!\). Thus for all \(u\in\R\),
\[
 |\cos(su)-L_s(u)|
     \le \frac{|s|^{2N}}{(2N)!}P_N(u)^2.
\]
At nodes both sides are zero. Polynomial exactness and
interpolation give a quadrature error for \(\cos(su)\) of
absolute value at most \(h_N^*|s|^{2N}/(2N)!\).

For a standard normal \(G\), set \(s=x+\sqrt{2\tau}G\).
Then \(\mathbb E[\cos(su)]=\ee^{-\tau u^2}\cos(xu)\).
Fubini applies by boundedness of the real cosine and finiteness
of both measures. They are even, so the same cosine gives the
exact and finite scalar readouts. The averaged error is bounded
by \(h_N^*\mathbb E[s^{2N}]/(2N)!\). Expanding with odd
Gaussian moments zero and
\(\mathbb E[G^{2k}]=(2k)!/(2^kk!)\) gives the displayed
sum. The Gaussian moments follow by integration by parts,
starting from \(\mathbb E[1]=1\).
\end{proof}

The corner \((\tau,x)=(0,0)\) has zero scalar error. At \(x=0\)
the allowance is \(h_N^*\tau^N/N!\) in base-mass units. At
\(\tau=0\) it is \(h_N^*x^{2N}/(2N)!\). There is also a
source-moment bound which uses only evenness and a Gaussian
moment of the base measure. Put \(M_1=\int\ee^{u^2}\dd\mu_*\).
The monic polynomial minimizes the squared norm among monic
degree-\(N\) polynomials, so
\[
 h_N^*\le\int u^{2N}\dd\mu_*\le N!M_1 .
\]
On \(0\le\tau\le1/2,\ |x|\le X\), reindexing the sum by
\(j=N-k\) gives
\[
 N!S_N(\tau,x)\le 2^{-N}\sum_{j=0}^N
   \frac{(2NX^2)^j}{(2j)!}\le 2^{-N}\cosh(X\sqrt{2N}).
\]
This uses \(N!/(N-j)!\le N^j\) and proves convergence on each
fixed real rectangle even without a coefficient envelope.
For a normalized tolerance \(\epsilon>0\), let
\(\ell=\log2\), \(b=\sqrt2 X\), and
\(d=\max(0,\log(M_1/\epsilon))\). Using \(\cosh t\le\ee^t\)
for \(t\ge0\), it suffices to take
\[
 N\ge\max\left(1,\left\lceil
     \left(\frac{b+\sqrt{b^2+4\ell d}}{2\ell}\right)^2
   \right\rceil\right).
\]
This is the nonnegative root test for
\(\ell(\sqrt N)^2-b\sqrt N-d\ge0\).

\section{Complete-source depth bound}\label{sec:depth}

We derive the logarithmic monic-norm estimate used for the
scalar remainder. For \(0\le\lambda\le1/2\), let \(P_n(u;\lambda)\)
be monic orthogonal for \(W_\lambda(u)\dd u\), and put
\[
 h_n(\lambda)=\int_{\R}P_n(u;\lambda)^2W_\lambda(u)\dd u,
 \qquad \beta_0=0,\qquad
 \beta_n(\lambda)=h_n(\lambda)/h_{n-1}(\lambda)\quad(n\ge1).
\]
The recurrence at each \(\lambda\) follows by the same parity
and orthogonality argument as at the base slice. Since multiplication
of the measure by a positive constant preserves the monic
polynomials, \(h_N^*=h_N(1/2)/h_0(1/2)\).

\begin{lemma}[Logarithmic coefficient envelope]\label{lem:log}
For every \(n\ge0\) and \(0\le\lambda\le1/2\),
\[
 \beta_n(\lambda)+\beta_{n+1}(\lambda)
       \le B_n,\qquad B_n=\frac{\arcosh^2(2n+12)}{16}.
\]
Consequently, with \(L_N=\arcosh(2N+10)\),
\begin{equation}\label{eq:log}
 h_N^*\le\prod_{j=0}^{N-1}B_j\le (L_N/4)^{2N},\qquad N\ge1.
\end{equation}
\end{lemma}
\begin{proof}
Write \(V_\lambda=-\log W_\lambda\), and on \(u\ge0\) use the
positive summands \(s_m=\ee^u r_m(2r_m-3)\ee^{-r_m}\).
Their logarithmic derivatives are
\[
 \frac{s_m'}{s_m}=9+\frac{12}{2r_m-3}-4r_m.
\]
The score \(\Phi'/\Phi\) is the positive weighted mean of these
derivatives by local uniform convergence. Using \(r_m\ge\pi
\ee^{4u}\), \(12/(2r_m-3)<4\), and \(2\lambda u\le u\) gives
\(V_\lambda'(u)\ge4\pi\ee^{4u}-13-u\) on \(u\ge0\).
Evenness now gives a useful all-line inequality,
\begin{equation}\label{eq:score}
 uV_\lambda'(u)\ge\cosh(4u)-11.
\end{equation}
Indeed, writing \(t=|u|\), for \(t\le1/2\) the lower bound
\(-13t-t^2\ge-27/4\) suffices, since \(\cosh 2<4\).
The series for \(\cosh 2\) after \(1+2\) starts with \(2/3\)
and the successive ratios are at most \(2/15\), which proves
that numerical inequality. For \(t\ge1/2\),
\(4\pi-1/t>10\) and \(10\ee^{4t}\ge10+40t\ge13+t\).
Thus \((4\pi t-1)\ee^{4t}\ge t(13+t)\), yielding
\(uV_\lambda'(u)\ge\ee^{4t}\ge\cosh(4t)\).

Fix \(n,\lambda\) and use the probability measure
\(\dd\nu_n=P_n(u;\lambda)^2W_\lambda(u)\dd u/h_n(\lambda)\).
Integrate \((uP_n^2 W_\lambda)'\) over the line. The bounds for
\(\Phi,\Phi'\) in Lemma~\ref{lem:domination} justify all integrals
and kill the boundary term. The polynomial \(uP_n'-nP_n\)
has degree less than \(n\) (and vanishes for \(n=0\)), so
orthogonality gives \(\mathbb E_{\nu_n}[uV_\lambda']=2n+1\).
Squaring the monic recurrence and integrating gives
\(\mathbb E_{\nu_n}[u^2]=\beta_n+\beta_{n+1}\), including
\(n=0\) by omitting the backward edge. The series of
\(v\mapsto\cosh(4\sqrt v)\) has nonnegative coefficients in
integer powers of \(v\ge0\) and is convex. Jensen and
\eqref{eq:score} give
\[
 \cosh(4\sqrt{\beta_n+\beta_{n+1}})
      \le\mathbb E_{\nu_n}[\cosh(4u)]\le 2n+12 .
\]
Taking the nonnegative inverse proves the first bound. At
\(\lambda=1/2\), use \(\beta_j\le B_{j-1}\) and telescope the
ratios \(h_j/h_{j-1}\). Monotonicity of \(B_j\) gives
\eqref{eq:log}.
\end{proof}

\begin{theorem}[Uniform real scalar depth]\label{thm:depth}
For \(N\ge1\) and \(X\ge0\),
\begin{equation}\label{eq:uniform}
 \sup_{\substack{0\le\tau\le 1/2\\ |x|\le X}}
 \frac{|H_{1/2-\tau}(x)-H_N(\tau,x)|}{H_{1/2}(0)}
 \le \exp\left(-N\log4+\frac{XL_N}{2}+\frac{L_N^2}{8}\right).
\end{equation}
If \(0<\epsilon\le1\) and \(K=X+\log(1/\epsilon)+1\), the
right side is at most \(\epsilon\) for
\begin{equation}\label{eq:depth}
 N=\left\lceil4K\log(16K)\right\rceil .
\end{equation}
\end{theorem}
\begin{proof}
The coefficient of \(r^{2N}\) in \(\ee^{\tau r^2}\cosh(xr)\)
is \(S_N(\tau,x)\), and every coefficient is nonnegative.
For \(r>0\),
\[
 S_N(\tau,x)\le r^{-2N}\ee^{\tau r^2+Xr}.
\]
Choose \(r=L_N/2\) and combine \eqref{eq:hermite} with
\eqref{eq:log} to obtain \eqref{eq:uniform}. For
\eqref{eq:depth} set \(A=\log(16K)\) and \(D=\log(1/\epsilon)\).
Since \(K\ge1\), \(\log K\le K-1\) and \(\log16<3\) give
\(A\le3K\). The ceiling satisfies \(4KA\le N<4KA+1\), so
\[
 4N+20\le16KA+24\le48K^2+24\le256K^2,\qquad
 L_N\le\log(4N+20)\le 2A .
\]
Using \(\log4>1\), the exponent in \eqref{eq:uniform}
is at most
\[
 -4KA+XA+A^2/2\le-\tfrac32 KA\le-D,
\]
since \(X\le K,\ D\le K,\ A>1\). This proves the claimed
tolerance for the stated integer \(N\), including \(X=0\)
or \(\epsilon=1\).
\end{proof}

The depth is sufficient for the scalar error with exact
coefficients and fixed mass normalization. Evaluation uses the
\(N-1\) internal edges of \(J_N\) and \(Z_*\). A sharper allowance
using the exact \(h_N^*\) can also require the omitted edge \(a_N\).
Both ways of evaluating the matched \(N\)-node rule have the same
output by Proposition~\ref{prop:gauss}.

\section{Finite coefficient uncertainty}\label{sec:uncertainty}

\begin{proposition}[Perturbation of the finite readout]\label{prop:perturb}
Let \(\widehat J\) be a real symmetric zero-diagonal Jacobi
matrix of size \(N\) sharing the tridiagonal pattern of \(J_N\).
Let \(\widehat Z\) be the midpoint of a certified positive interval
for \(Z_*\), and let \(\eta_Z\) be its radius divided by its lower
endpoint. Define the ideal nominal physical scalar by
\[
 \widehat H(\tau,x)=\frac{\widehat Z}{2}
   \ip{e_0}{\ee^{-\tau\widehat J^2}\cos(x\widehat J)e_0}.
\]
If \(\delta\ge\|J_N-\widehat J\|_2\), then on the real domain
of Theorem~\ref{thm:hermite},
\begin{equation}\label{eq:perturb}
 \frac{|H_{1/2-\tau}(x)-\widehat H(\tau,x)|}{Z_*/2}
 \le h_N^*S_N(\tau,x)
       +\left(|x|+\sqrt{4\tau/\pi}\right)\delta+\eta_Z.
\end{equation}
When each nominal internal edge is the exact midpoint of an
enclosure of the corresponding exact edge, one may take \(\delta\)
to be the maximum sum of the radii of the two edges incident to
a site, with absent edge radii set to zero.
\end{proposition}
\begin{proof}
For real symmetric finite matrices \(A,B\), differentiate
\(\ee^{\ii(s-v)A}\ee^{\ii vB}\) between \(v=0\) and \(v=s\).
The resulting integral and unitarity give
\(\|\ee^{\ii sA}-\ee^{\ii sB}\|_2\le |s|\|A-B\|_2\),
also for \(s<0\) with reversed integration limits. Averaging
the two complex exponentials yields the same bound for the
difference of cosines. Apply the Gaussian average of
Theorem~\ref{thm:hermite} through the finite spectral calculus.
It replaces \(\cos(sA)\) by \(\ee^{-\tau A^2}\cos(xA)\).
The bound \(\mathbb E|s|\le|x|+\sqrt{4\tau/\pi}\) follows
from \(\mathbb E|G|=\sqrt{2/\pi}\). Parity makes the original
readout real for both matrices, and the nominal scalar matrix
element in \eqref{eq:perturb} has modulus at most \(1\).
Splitting the scalar difference at the exact mass and the
nominal matrix costs at most \(\eta_Z\) for the mass and the
displayed perturbation for the coefficients. Combine this
with \eqref{eq:hermite}. The symmetric difference matrix has
operator norm bounded by its maximum absolute row sum,
which only uses the two internal edges adjoining each row.
\end{proof}

\subsection{A 16-site interval example}\label{sec:example}

An existing finite acquisition \cite{certificate} encloses the
17 even raw moments of \(W_*\dd u\) through degree \(32\).
The moment integrals restore both the omitted arithmetic modes
and the spatial tails. In the real monomial Gram matrix, entries
of opposite parity vanish. Eliminating each parity block by
interval \(LDL^{\mathsf T}\) encloses the squared monic norms
\(h_n(1/2)\), since the pivot is the squared norm of the monomial
after projection off lower degrees of its parity. Positive
pivots permit the successive square roots of their ratios,
enclosing \(a_1,\ldots,a_{16}\). The ideal nominal 16-site
matrix uses the first 15 midpoints; the last coupling belongs
to the error norm. The ratio \(h_{16}(1/2)/h_0(1/2)\) directly
pays the truncation with \eqref{eq:hermite}.

On \(0\le\tau\le1/2,\ |x|\le8\), the matrix factor in
\eqref{eq:perturb} is below \(9\). The certificate records the
following upward bounds at 256 bits, in units of \(H_{1/2}(0)\).
\begin{center}
 \begin{tabular}{@{}lr@{}}
 \toprule
 Component & Upper bound\\
 \midrule
 Uniform truncation \(h_{16}^*S_{16}(1/2,8)\)
       & \(1.515\times10^{-25}\)\\
 Uniform coefficient uncertainty \(9\delta\)
       & \(1.385\times10^{-46}\)\\
 Uniform mass uncertainty \(\eta_Z\)
       & \(3.472\times10^{-68}\)\\
 Largest nominal rounding allowance at the nine samples
       & \(1.740\times10^{-77}\)\\
 \bottomrule
 \end{tabular}
\end{center}
The exact enclosures before upward decimal display certify
total normalized absolute error below \(1.515\times 10^{-25}\)
for the ideal nominal scalar on the entire rectangle. The
same bound includes rounding at the nine evaluated points
\(\tau\in\{0,1/4,1/2\}\), \(x\in\{0,4,8\}\).
The 128-bit acquisition had positive pivots and couplings but
missed its coefficient-width allowance, with \(9\delta\) about
\(1.791\times10^{-7}\) against its \(2\times10^{-9}\) cap.
The successful 256-bit acquisition also encloses the infinity-norm
condition numbers of its diagonally scaled even and odd Gram
blocks by \(6.94\times10^7\) and \(2.02\times10^7\), respectively.
These are diagnostics of this finite procedure and its data.

\section{Outstanding challenges}\label{sec:challenges}

\paragraph{State transport.}
The scalar approximation leaves open a quantitative bound for
the norm outside the first \(N\) polynomial-degree sites of the
exact state. Section~\ref{sec:mps} proves escape from every fixed
prefix while leaving open the rate and the dominant degrees.
An estimate for the projected state would address a different
norm from \eqref{eq:hermite}. At the operator level any real
nonzero height increment \(d\) satisfies
\(\|\ee^{\ii dJ_*}-I\|=2\), since the essential supremum of
\(|\ee^{\ii du}-1|\) is attained as a supremum on the full
support. The height group is nevertheless strongly continuous
by dominated convergence in the multiplier representation.
For the prepared orbit, \(|\ee^{\ii du}-1|\le|du|\) even gives
\(\|v(\tau,x+d)-v(\tau,x)\|\le|d|\|J_*\ee^{-\tau J_*^2}e_0\|\),
uniformly in \(x\).

\paragraph{Growing height and conditioning.}
The sufficient depth in \eqref{eq:depth} follows from an upper
envelope. Sharper estimates in a joint limit of height and
accuracy would require quantitative information beyond that
used here. The interval example leaves open how Gram conditioning
and required coefficient precision scale with \(N\), how best
to enclose the coefficients from the complete source, and what
work repeated evaluations require after acquisition. For the
same measure and node count, direct Gaussian quadrature uses
the identical scalar rule.

\paragraph{Signed readout.}
The error scale is the fixed positive mass \(H_{1/2}(0)\).
A relative bound for a varying readout would also need a
nonzero lower margin. At any zero in a covered rectangle an
absolute allowance necessarily dominates the absolute approximate
value. For example, \(2+2\ee^{\lambda-T}\cos x\), with
\(0<T\le1/2\), comes from a positive even measure on
\(\{-1,0,1\}\) at the base slice and satisfies the same heat
equation. The full three-site rule is exact there, while the
field has a double zero at Newman time \(\lambda=T\) and
height \(x=\pi\). The example terminates
at finite support and has no full-line score. Even continuous
modifications built from positive weighted square summands can
retain the score proof in Lemma~\ref{lem:log} and allow
collisions \cite{controls}. Their locally absolutely continuous
even extensions suffice for its integration by parts. A signed
estimate for the complete theta field must supply further
information sensitive to cancellations in the physical cosine
readout.

\section*{Acknowledgments}

Draft prepared with substantial OpenAI Codex LLM assistance.
The execution dispatch selected GPT-6.1 Sol with Ultra effort.
The backend model and runtime effort were unavailable for
independent verification. The internal assessment shares the
draft's sources and constitutes an author-side check.

\begin{thebibliography}{9}

\bibitem{theta}
NIST Digital Library of Mathematical Functions.
\emph{Theta Functions}, \S20.7(viii), equation 20.7.32.
\url{https://dlmf.nist.gov/20.7.E32}.

\bibitem{chin}
A.~W.~Chin, A.~Rivas, S.~F.~Huelga, and M.~B.~Plenio.
Exact mapping between system-reservoir quantum models and
semi-infinite discrete chains using orthogonal polynomials.
\emph{J. Math. Phys.} \textbf{51}, 092109 (2010).
\href{https://doi.org/10.1063/1.3490188}{doi:10.1063/1.3490188}.
\href{https://arxiv.org/abs/1006.4507v2}{arXiv:1006.4507v2}.

\bibitem{gauss}
NIST Digital Library of Mathematical Functions.
\emph{Numerical Methods}, \S3.5(v)--(vi).
\url{https://dlmf.nist.gov/3.5}.

\bibitem{certificate}
E.~Baker, with OpenAI Codex assistance.
\emph{Certified complete-source acquisition and readout}.
Shifted Zeta Positivity, Investigation 6, 14 November 2026.
Execution record and exact-rational companion certificate.

\bibitem{controls}
Shifted Zeta Positivity project.
\emph{Collision controls and the screening they impose}.
Result summary 12, with catalogued weighted-square examples and
their source certificates.

\end{thebibliography}
\end{document}
